\chapter{A Worked C$--$ Translation Example}
\label{ch:workcmm}

\section{Setup}
\label{sec:work-setup}

To make the middle of the pipeline concrete, we follow one artifact ---
the summation of $v = (1,2,3,4)$ --- from the IR through the C$--$
emission and the static/FOL/ASP checks that precede the bridge. This is
the same artifact traced at the APL level in \cref{ch:apl-sem}.

\section{IR In}
\label{sec:work-ir}

The Layer~2 IR for this artifact (after the contract of
\cref{ch:apl-sem} holds) is the decorated record:
\begin{center}
\texttt{\{ shape=[4], dtype=int, expr=SUM(v), cert=lean:arith\_sum,
  receipt=r1 \}}
\end{center}

\section{C$--$ Out}
\label{sec:work-out}

Layer~5 (Liquid Haskell $\to$ C$--$) emits the following
verified-translation unit:
\begin{lstlisting}[language=C,caption={C$--$ emitted for $\Sigma\,\omega$
  on $v=(1,2,3,4)$.},label=lst:work-cmm]
section "DATA" {
  v_arr:
    bits64[4] = {1, 2, 3, 4};
}
f_sumV :: bits64
f_sumV {
  i0: bits64, r0: bits64;
  i0 = 0;
  r0 = 0;
  loop:
    if (i0 >= 4) goto done;
    r0 = r0 + v_arr[i0];
    i0 = i0 + 1;
    goto loop;
  done:
    return (r0);
}
\end{lstlisting}
The loop is a left fold over the array; integer addition is associative,
so it computes the same value as the APL reduction $+/\,v$
(\cref{ch:apl-sem}).

\section{Static Analysis}
\label{sec:work-static}

Layer~6 confirms there are no out-of-bounds reads: $i0$ is incremented
only after a guard \texttt{i0 >= 4}, and the array has length 4. It also
proves \texttt{r0} is initialized before \texttt{return}, removing the
undefined-read class of bugs. The produced invariant is
\[
0 \le i0 \le 4 \;\wedge\; r0 = \sum_{k=0}^{i0-1} v_k.
\]

\section{First-Order Logic Emission}
\label{sec:work-fol}

Layer~7 renders the same property as a first-order clause for the external
prover:
\[
\forall i_0,r_0.\; \bigl(\text{loop}(i_0,r_0) \to
(i_0\le 4 \wedge r_0 = \text{sum}(v,i_0))\bigr).
\]
The FOL form is what the ASP solver consumes, decoupling the proof search
from the C$--$ syntax.

\section{ASP Schedule}
\label{sec:work-asp}

Layer~8 packages the clause as an answer-set program with the goal
\texttt{:- not holds(f\_sumV = 10).} The solver finds a stable model in
which \texttt{holds(f\_sumV = 10)} is true, matching the closed-form
expectation $10$. The schedule is:
\begin{enumerate}
  \item Encode invariant as facts.
  \item Assert goal negation as a constraint.
  \item Search for a model; if found, the constraint is violated, proving
        the goal.
\end{enumerate}

\section{Receipt Chain}
\label{sec:work-receipt}

Each layer appends a receipt referencing the previous one:
\[
r_1 \to r_2 \to r_3 \to r_4 \to r_5 \to r_6 \to r_7 \to r_8,
\]
and the chain is sealed into the WORM ledger (\cref{ch:worm}). A verifier
recomputing \texttt{SHA-256} of each artifact and checking the Ed25519
Plasma-Gate signatures obtains a complete proof that the C$--$ function
\texttt{f\_sumV} returns $10$ for the given input, with no step trusted by
assertion alone.

\section{Why This Matters}
\label{sec:work-why}

The example is trivial arithmetically, but the \emph{structure} of the
proof is the product: every non-trivial summation in the system is
subject to the identical, machine-checked chain. Scaling the input from
four elements to four million changes only the runtime of the loop, not
the force of the certificate.
